\section{Amplification to uniform property \texorpdfstring{$\Gamma$}{Gamma}}\label{n:sec:amplification}

The small splittings from Section~\ref{n:sec:small} have errors of order
$\lambda^{3/2}$. At a fixed target proportion, a number of steps
proportional to $\lambda^{-1}$ therefore accumulates an error of order
$\lambda^{1/2}$. We transfer the small splittings to arbitrary separable
coefficient algebras and multiply commuting complements. This gives any
fixed trace proportion with vanishing error, after which a bounded
diagonal argument gives exact projections.

\begin{lemma}\label{n:lem:relative-small}
For every separable unital $E\subset D$ and $0<\lambda<1$, there is a
positive contraction $e\in D\cap E'$ such that, for $\tau\in \Lambda$,
\begin{equation}\label{n:eq:relative-small}
 |\tau(ex)-\lambda\tau(x)|\le\lambda^{3/2}
 \quad(x\in E_{\mathrm{sa}},\ \|x\|\le1),
 \qquad \tau(e-e^2)\le\lambda^{3/2}.
\end{equation}
\end{lemma}
\begin{proof}
Choose a dense sequence of self-adjoint contractions in $E$ and represent
each by self-adjoint contractions in $A$. At coordinate $n$, apply
Proposition~\ref{n:prop:small} to the first $n$ entries of these
representatives, with $\rho=1/n$. The resulting sequence gives a positive
contraction in $D$ whose commutators with the dense sequence have zero
uniform $2$-norm. They are therefore zero in $D$. Passing to the limit
traces proves the inequalities on the dense sequence, and norm continuity
extends them to the entire self-adjoint unit ball.
\end{proof}

\begin{proposition}\label{n:prop:split}
For every separable unital $E\subset D$ and $0<c<1$, there is a projection
$f\in D\cap E'$ such that
\begin{equation}\label{n:eq:split}
 \tau(fx)=c\tau(x)\qquad(x\in E,\ \tau\in \Lambda).
\end{equation}
\end{proposition}
\begin{proof}
For an integer $m\ge1$, put $\lambda=1-c^{1/m}$. Successively choose
$e_1,\ldots,e_m$ from Lemma~\ref{n:lem:relative-small}, enlarging $E$ to
contain the previous choices at each step. They commute, and
\[
 P_j=\prod_{i=1}^j(1-e_i),\qquad P_0=1,
\]
are positive contractions commuting with $E$. For a self-adjoint
contraction $x\in E$, the element $P_{j-1}x$ is a permitted test at step
$j$. The resulting recurrence gives
\[
 |\tau(P_mx)-c\tau(x)|\le m\lambda^{3/2}.
\]
For commuting positive contractions $a,b$,
\[
 ab-(ab)^2=b(a-a^2)+a^2(b-b^2)\le(a-a^2)+(b-b^2).
\]
Thus $\tau(P_m-P_m^2)\le m\lambda^{3/2}$ as well. Since
$\lambda\le(-\log c)/m$, both errors tend to zero.

We spell out the diagonal passage because $E$ need not consist of constant
sequences. Fix representatives for a dense sequence of self-adjoint
contractions of $E$. For each $k$, choose $m$ large enough and a positive
contraction lift $P_n^{(k)}$ so that on an $\omega$-large set $G_k$, the
first $k$ commutators and trace moments, and the uniform $2$-norm projection
defect, are all within $1/k$ of their targets. Such a set exists by the
uniform limit-trace estimates and exact commutation in $D$; use
$\tau((P_m-P_m^2)^2)\le\tau(P_m-P_m^2)$ for the defect. Put
\[
 F_k=\{n:n\ge k\}\cap G_1\cap\cdots\cap G_k.
\]
These sets decrease, belong to $\omega$, and have empty intersection. At
coordinate $n\in F_1$, use $P_n^{(k(n))}$ for the largest $k(n)$ with
$n\in F_{k(n)}$; use zero off $F_1$. Then $k(n)\to_\omega\infty$.
The resulting class is a projection commuting with $E$ and satisfying
\eqref{n:eq:split} on the dense sequence. Continuity and complex linearity
complete the proof.
\end{proof}

\begin{theorem}\label{n:thm:gamma}
For every $n\ge2$ there are orthogonal projections $f_1,\ldots,f_n\in
D\cap A'$ summing to $1_D$ such that
\[
 \tau(f_i a)=\frac1n\tau(a)\qquad(a\in A,\ \tau\in \Lambda).
\]
In particular, $A$ has uniform property $\Gamma$.
\end{theorem}
\begin{proof}
Start with the separable unital algebra $C^*(A,1_D)$. At step
$j=1,\ldots,n-1$, adjoin the previously selected projections, apply
Proposition~\ref{n:prop:split} with $c=1/(n-j+1)$, and multiply its
projection by the remaining complement $q$. The factors commute, so this
gives a projection $f_j\le q$. Apply \eqref{n:eq:split} to $qa$ in the
enlarged algebra. Inductively $\tau(qa)=(n-j+1)\tau(a)/n$, so
$\tau(f_ja)=\tau(a)/n$. Take the final complement for $f_n$.
These are exactly the projection identities in the definition of uniform
property $\Gamma$ \cite[Definition 2.1]{CETW}.
\end{proof}
